%-------------------------------------------------------------------------------
%	SECTION TITLE
%-------------------------------------------------------------------------------
\cvsection{Side Projects}


%-------------------------------------------------------------------------------
%	CONTENT
%-------------------------------------------------------------------------------
\begin{cventries}

%---------------------------------------------------------
  \cventry
    {C / C++ / Assembly (ARMv8, RISC-V)} % Affiliation/role
    {CS Self-Learning} % Organization/group
    {Open Courseware} % Location
    {Apr. 2022 - Present} % Date(s)
    {
      \begin{cvitems} % Description(s) of experience/contributions/knowledge
        \item{Implemented small kernels on ARMv8 (Raspberry Pi) and RISC-V (QEMU), covering processor initialization, memory management, interrupt handling, context switching, process scheduling, and system calls.}
        \item{Extended NachOS with page tables for multi-programming, I/O system calls, and a multilevel feedback queue scheduler with aging.}
        \item{Built a Linux kernel module (/dev/fibonacci) with big-number arithmetic, \textbf{improving performance by \textasciitilde75\%} via algorithm optimization.}
        \item{Currently taking NJU's 2026 \textbf{AI-Native Software Engineering} course, covering specs, CI, and testing for agent-driven development.}
        \item {Site: \href{https://tingggggg.github.io/cs-self-learning/}{\cvlink{CS Self-Learning}} $\cdot$ Code: \href{https://github.com/tingggggg/NTHU_Nachos_Project}{\cvlink{NachOS}} $\cdot$ \href{https://github.com/tingggggg/OSDIg}{\cvlink{Raspberry Pi}} $\cdot$ \href{https://github.com/tingggggg/riscv-operating-system/tree/main/os}{\cvlink{RISC-V}} $\cdot$ \href{https://github.com/tingggggg/fibdrv}{\cvlink{fibdrv}} $\cdot$ Notes: \href{https://github.com/tingggggg/NJU_AI-Native_Software_Engineering_2026}{\cvlink{AI-Native SE}}}
      \end{cvitems}
    }

%---------------------------------------------------------
\end{cventries}
